\section{Introduction}
\label{introduction}

Natural language requirements are often written around intended functionality, while the environment in which the software will be deployed remains implicit.
Missing information about devices, runtime conditions, user situations, and deployment assumptions can lead to incomplete or misleading requirements.
This paper frames environment-aware requirement elicitation as a research problem in intelligent requirements engineering.

\subsection{Background and Motivation}
\label{intro:background}

Environment-aware requirement engineering aims to capture not only what a system should do, but also where, under which conditions, and for whom it should operate.
This subsection motivates the need to elicit environmental context early, before requirements are refined into downstream artifacts.
Recent work has also emphasized requirements as a central interface for future software engineering workflows~\cite{robinsonRequirementsAreAll2025}.

\subsection{Challenges of Environment-Aware Requirement Elicitation}
\label{intro:challenges}

The paper focuses on four recurring challenges:
\begin{itemize}[leftmargin=*]
  \item natural language requirements often omit deployment and operational assumptions;
  \item stakeholders may describe user goals without specifying devices, locations, connectivity, or runtime constraints;
  \item static requirement generation cannot reliably infer environment-specific needs;
  \item elicitation tools need to ask targeted clarification questions without overburdening human stakeholders.
\end{itemize}

\subsection{Multi-Agent Context Understanding}
\label{intro:multi_agent_context}

We investigate whether a lightweight multi-agent workflow can improve context understanding during elicitation.
The key idea is to separate stakeholder interviewing, end-user perspective simulation, and environment/deployment analysis while keeping the workflow centered on missing environmental information.

\subsection{Contributions}
\label{intro:contributions}

The intended contributions are:
\begin{itemize}[leftmargin=*]
  \item a research problem formulation for environment-aware requirement elicitation;
  \item a multi-agent workflow for detecting missing environmental context;
  \item a knowledge-assisted context understanding method for scenario refinement;
  \item an empirical evaluation plan based on context completeness, environment consistency, requirement completeness, and clarification effectiveness.
\end{itemize}
